\documentclass[twocolumn]{iopjournal}

\usepackage{amsmath,amssymb,amsfonts}

\fancyhfoffset[L]{0mm}
\setlength{\marginparwidth}{0pt}
\setlength{\marginparsep}{0pt}

% ==========================================
% ZENODO PREPRINT HEADERS (ACTIVE)
% ==========================================
\renewcommand{\articletype}[1]{%
  {\vspace*{-8mm}\noindent \Large \sf Preprint}\par
  \vspace*{6mm}\noindent{\scriptsize \sf{\bfseries \MakeUppercase{#1}}}}

\fancyhead[L]{}
\fancyhead[R]{}

\begin{document}

\articletype{Supplementary Material}

\title{Supplementary Material: The Constitutive Law\\
Read from the Metric}

\author{Stephen Nelson}

\affil{Field Dynamics Research}

\ead{stephen@fielddynamics.org}

\begin{abstract}
This note derives the gravitational constitutive law from the
metric and the tetrad equation through a sequence of four
operations: reading numbers from
$\eta_{ab} = \mathrm{diag}(-1,+1,+1,+1)$, noting the degree
of the tetrad equation, performing arithmetic, and applying
algebra. The metric is $4 \times 4$ ($K = 4$) with three
positive eigenvalues ($D = 3$). The tetrad equation is degree~2
(two roots). These numbers satisfy $2(D+1) = 2^D$; $D = 3$ is
the minimum spatial dimensionality for which this closure relation
can hold, and in fact its only positive integer solution. The tetrad equation contains
three structural ingredients whose graded dimension is
$K^0 + K^1 + K^2 = 1 + 4 + 16 = 21$. Parity interpolation
on the three states makes the polynomial linear in its state
parameter. Normalising the resulting fractional linear
composition by the polynomial's range yields
$\mu(x) = x/(x+1)$. Two companion supplements derive the same
function by independent methods: one from equilibrium field
closure, the other from boundary conditions on the quadratic
tetrad equation. The convergence of three routes with no shared
method on a single function is the subject of the main letter.
\end{abstract}

% ----------------------------------------------------------------------

\section*{Read}

The internal metric of the tetrad formalism is
\begin{equation}
  \eta_{ab} = \mathrm{diag}(-1,+1,+1,+1).
  \label{eq:eta}
\end{equation}
Three numbers are read directly from this object:

\paragraph{$K = 4$.} The matrix is $4 \times 4$. The frame
index $a$ runs over four values. This is the frame dimension.

\paragraph{$D = 3$.} The matrix has three positive eigenvalues
and one negative. The count of positive eigenvalues is the
spatial dimension.

\paragraph{$n = 2$.} The tetrad equation
$g_{\mu\nu} = \eta_{ab}\, e^a{}_\mu\, e^b{}_\nu$ contains $e$
twice on the right-hand side: it is degree~2 in $e$. A degree-2
equation has two roots. In the spatial sector, the two roots are
parity-related tetrahedra~\cite{Nelson2026scheme}. The count
$n = 2$ is the degree of the equation, visible by inspection.

$K = 4$ and $D = 3$ are read directly from the metric. The
count $n = 2$ is read from the quadratic degree of the tetrad
equation and geometrically realised as the two parity-related
bounded regions established in the supplementary
construction~\cite{Nelson2026scheme}.

% ----------------------------------------------------------------------

\section*{Arithmetic}

From $K = 4$, $D = 3$, and $n = 2$, form two expressions:

\begin{align}
  \text{Left:} \quad n(D + 1) &= 2(3 + 1) = 8, \label{eq:left} \\
  \text{Right:} \quad n^D &= 2^3 = 8. \label{eq:right}
\end{align}

They are equal. This is arithmetic, not a postulate.

The left side counts faces: two roots, each carrying a $D$-simplex
with $D + 1 = K$ faces. The right side counts vertices: $D$ binary
axes, each admitting two orientations, giving $2^D$ sign
combinations. The main letter~\cite{Nelson2026geometry} proved
that these are the half-space intersections of the four derived
directions, and that the tetrad equation produces exactly two
bounded tetrahedral regions, TetA and TetB, each with $K = 4$
faces. Ghost-free bimetric gravity likewise requires exactly two
dynamical tetrad fields, each with $K = 4$ frame
legs~\cite{HR2012a}. The supplementary isotropy
proof~\cite{Nelson2026signature} establishes that the four spatial
legs of the tetrad are the four face normals of a regular
tetrahedron: frame legs and tetrahedral faces are the same four
directions. The same counts therefore appear in both
contexts: $2K = 8$ as the doubled frame count, and $K^2 = 16$ as
the number of frame-index pairings $(a,b)$ in the $4 \times 4$
coupling structure. Their recurrence across the metric
reading, the derived geometry, and the published bimetric framework
confirms that these integers are not arbitrary.

The equality $2(D+1) = 2^D$ is a Diophantine equation in $D$.
Checking whether any other positive integer satisfies it:

\vspace{4pt}
\noindent
\begin{tabular*}{\columnwidth}{@{\extracolsep{\fill}}cccc@{}}
\hline
$D$ & $2(D+1)$ & $2^D$ & Equal? \\
\hline
1 & 4 & 2 & No \\
2 & 6 & 4 & No \\
\textbf{3} & \textbf{8} & \textbf{8} & \textbf{Yes} \\
4 & 10 & 16 & No \\
5 & 12 & 32 & No \\
\hline
\end{tabular*}
\vspace{4pt}

\noindent $D = 3$ is the minimum spatial dimensionality at which
the closure relation can hold: $D = 1$ and $D = 2$ are
insufficient, $D = 3$ succeeds, and the exponential side
thereafter outgrows the linear side strictly, so $D = 3$ is also
the unique positive integer solution.

% ----------------------------------------------------------------------

\section*{Verification}

The two integers $K = 4$ and $D = 3$ recover the principal
dimensional and combinatorial properties of $g_{\mu\nu}$ shown
below by arithmetic alone:

\vspace{4pt}
\noindent
\begin{tabular*}{\columnwidth}{@{\extracolsep{\fill}}lcc@{}}
\hline
Property & $K, D$ & Known \\
\hline
Matrix size & $K \times K = 4 \times 4$ & $g_{\mu\nu}$: $\mu,\nu = 0..3$ \\
Signature & $(-) + (+)^D$ & $(-,+,+,+)$ \\
Independent entries & $K(K\!+\!1)/2 = 10$ & Symmetric $4 \times 4$ \\
Total entries & $K^2 = 16$ & $4 \times 4$ \\
Spatial block & $D \times D = 3 \times 3$ & $g_{ij}$ \\
Frame legs & $K = 4$ & $e^a{}_\mu,\; a\!=\!0..3$ \\
Spacetime dimension & $D + 1 = 4$ & $3 + 1$ \\
Vertices of compound & $2^D = 8$ & Cube corners \\
Faces of compound & $2K = 8$ & Octahedron faces \\
f-vector sum & $K(K\!+\!D) = 28$ & $V + E + F$ \\
\hline
\end{tabular*}
\vspace{4pt}

\noindent Every entry matches. Two integers recover the
dimensional and combinatorial structure displayed above.

% ----------------------------------------------------------------------

\section*{The Polynomial}

The tetrad equation contains three structural ingredients on the
$K$-dimensional frame space. A scalar closure condition
$\det(e) \neq 0$ with no free frame index (rank~0, dimension
$K^0 = 1$). A frame $e^a{}_\mu$ with one free frame index
(rank~1, dimension $K^1 = K = 4$). A bilinear coupling
$\eta_{ab}\, e^a\, e^b$ with two free frame indices (rank~2,
dimension $K^2 = 16$). No rank-3 object appears in the equation
or its field equations~\cite{MTW}. Taking the graded dimension
of a rank-$r$ frame-space ingredient to be $K^r$, the natural
three-term encoding is:
\begin{equation}
  f(K) = K^0 + K^1 + K^2 = 1 + 4 + 16 = 21.
  \label{eq:poly}
\end{equation}
This is the finite geometric series $(K^3 - 1)/(K - 1) =
63/3 = 21$.

Under spatial parity $P: e^i \to -e^i$, a rank-$r$ tensor on the
frame space transforms with sign $(-1)^r$~\cite{MTW}. Rank~0
($(-1)^0 = +1$) and rank~2 ($(-1)^2 = +1$) are parity-even.
Rank~1 ($(-1)^1 = -1$) is parity-odd. Accordingly, the
parity-varying contribution is assigned to the rank-1 term:
\begin{equation}
  f(K, s) = 1 + sK + K^2.
  \label{eq:poly_s}
\end{equation}

Evaluating at the two roots ($s = \pm 1$) and their midpoint
($s = 0$):
\begin{equation}
  f(4,+1) = 21, \qquad f(4,0) = 17, \qquad f(4,-1) = 13.
\end{equation}

The range of the polynomial is:
\begin{equation}
  \Delta f = 21 - 13 = 2K = 8 = 2^D.
  \label{eq:range}
\end{equation}
This is the closure condition $2(D+1) = 2^D$ read as the
dynamic range of the polynomial. The same equation that makes
$D = 3$ unique also sets the normalisation denominator.

% ----------------------------------------------------------------------

\section*{The Constitutive Law}

The stellated octahedron has exactly three structural components,
each proved in the main letter and its
supplements~\cite{Nelson2026geometry,Nelson2026scheme}: TetA (one
bounded region, identity component of $O(3)$), TetB (the other
bounded region, parity component), and the octahedral core (the
intersection TetA~$\cap$~TetB, with six vertices at the
tetrahedral edge midpoints). These three components are the three
geometric states of the structure. The state parameter $s$ labels
them: $s = +1$ for TetA, $s = -1$ for TetB, $s = 0$ for the
octahedral core. Three components, three states. Not chosen;
counted.

The polynomial $f(K,s) = 1 + sK + K^2$ is linear in $s$.
This linearity follows from interpolation on three points with
odd symmetry: $s$ takes values $\{+1, 0, -1\}$, the parity-odd
requirement sets the $s$-dependent part to zero at $s = 0$, and
the endpoint values are $\pm K$. The unique polynomial of degree
$\leq 2$ passing through $(-1, -K)$, $(0, 0)$, $(+1, +K)$ is
$Ks$. Linear, by interpolation uniqueness, not by minimality
assumption.

The geometry has three components, so the fundamental theorem of
projective geometry applies: three distinct points on the
projective line uniquely determine a fractional linear map. Three
constraints fix the M\"{o}bius map $s(x)$
making the state parameter continuous in $x = g/a_0$: $s \to
+1$ as $x \to \infty$, $s = 0$ at $x = 1$ (a normalisation
convention defining $a_0$), and $s \to -1$ as $x \to 0$. Among
fractional linear maps satisfying these three conditions, the
unique solution is $s(x) = (x-1)/(x+1)$. The fractional linear
class is not selected but forced: the map $s(x)$ must be rational,
to preserve the polynomial's algebraic structure under composition,
and bijective, so that a single acceleration $x$ decodes to exactly
one geometric state $s$ and vice versa. A rational bijection on the
projective line is necessarily fractional linear; this is a
classification theorem of projective geometry, not a modelling
choice~\cite{Hartshorne1977}.
This forcing is confirmed by closure under inversion:
$\mu = x/(x+1)$ inverts to $x = \mu/(1-\mu)$, which is again
fractional linear, with no increase in algebraic degree. Transcendental or higher-degree rational
alternatives generally forfeit the same exact fractional-linear
closure under composition and inversion. This exact invertibility
is also what permits closed-form decoding from observed kinematics
to baryonic mass, rather than numerical root-finding.

Because $f(K,s)$ is linear in $s$ and $s(x)$ is fractional
linear, the composition $f(x)$ is fractional linear.
Normalising a fractional linear function preserves the form.
The constitutive law inherits its rational structure from the
polynomial's linearity in $s$.

Composing:
\begin{equation}
  f(x) = 1 + \frac{x - 1}{x + 1} \cdot 4 + 16
       = \frac{21x + 13}{x + 1}.
\end{equation}

Normalising by the range $\Delta f = 8$:
\begin{equation}
  \mu(x) = \frac{f(x) - 13}{8} = \frac{x}{x + 1}.
  \label{eq:mu}
\end{equation}

The ``1'' in the denominator corresponds to the rank-0 closure
term $K^0 = 1$, consistent with the metric's single timelike
direction. The normalisation denominator is
$2K = 2^D = 8$: the closure condition itself.

% ----------------------------------------------------------------------

\section*{The Complete Chain}

\begin{equation*}
  \boxed{\eta_{ab} = \mathrm{diag}(-1,+1,+1,+1)}
\end{equation*}
\vspace{-6pt}
\begin{equation*}
  \downarrow\; \textsc{read + degree}
\end{equation*}
\vspace{-6pt}
\begin{equation*}
  \boxed{K = 4, \quad D = 3, \quad n = 2}
\end{equation*}
\vspace{-6pt}
\begin{equation*}
  \downarrow\; \textsc{arithmetic}
\end{equation*}
\vspace{-6pt}
\begin{equation*}
  \boxed{2(3+1) = 2^3 = 8 \quad \checkmark \quad
  \text{threshold and unique solution}}
\end{equation*}
\vspace{-6pt}
\begin{equation*}
  \downarrow\; \textsc{arithmetic}
\end{equation*}
\vspace{-6pt}
\begin{equation*}
  \boxed{f = 1 + 4 + 16 = 21}
\end{equation*}
\vspace{-6pt}
\begin{equation*}
  \downarrow\; \textsc{parity} \to \text{interpolation} \to
  \text{linear in } s \to \textsc{m\"{o}bius} \to
  \textsc{normalise}
\end{equation*}
\vspace{-6pt}
\begin{equation*}
  \boxed{\mu(x) = \frac{x}{x+1}}
\end{equation*}

\vspace{6pt}
Every step is either direct reading, arithmetic,
representation-theoretic classification, or algebraic
normalisation. No empirical fit enters the derivation. The
fractional linear form of the constitutive law is inherited
from the polynomial's linearity in $s$, which is itself fixed
by interpolation uniqueness on the three parity states.

The function $\mu(x) = x/(x+1)$ was identified empirically by
Milgrom~\cite{Milgrom1983} and selected by fitting by
Famaey and Binney~\cite{FamaeyBinney2005}. The present note
obtains it from metric data and the quadratic tetrad structure
through graded rank counting, parity interpolation, and rational
normalisation. Two companion supplements obtain the same
function by independent methods: one from equilibrium field
closure, the other from boundary conditions on the quadratic
tetrad equation. The main letter presents all three routes and
their convergence.

% ----------------------------------------------------------------------

\begin{thebibliography}{99}

\bibitem{Nelson2026geometry}
S.~Nelson,
\textit{The Coordinate-Space Geometry of the Tetrad Frame},
Class.\ Quantum Grav.\ (2026), submitted.

\bibitem{Nelson2026scheme}
S.~Nelson,
\textit{Supplementary Material: The Tetrad Scheme and the Two
Bounded Regions} (2026).

\bibitem{Nelson2026signature}
S.~Nelson,
\textit{Supplementary Material: From the Lorentz Signature to the
Tetrahedral Frame} (2026).

\bibitem{MTW}
C.~W.~Misner, K.~S.~Thorne and J.~A.~Wheeler,
\textit{Gravitation}
(W.~H.~Freeman, San Francisco, 1973).

\bibitem{Hartshorne1977}
R.~Hartshorne,
\textit{Algebraic Geometry}
(Springer, New York, 1977).

\bibitem{Milgrom1983}
M.~Milgrom,
\textit{A modification of the Newtonian dynamics as a possible
alternative to the hidden mass hypothesis},
Astrophys.\ J.\ \textbf{270} (1983) 365.

\bibitem{FamaeyBinney2005}
B.~Famaey and J.~Binney,
\textit{Modified Newtonian dynamics in the Milky Way},
Mon.\ Not.\ R.\ Astron.\ Soc.\ \textbf{363} (2005) 603.

\bibitem{HR2012a}
S.~F.~Hassan and R.~A.~Rosen,
\textit{Bimetric gravity from ghost-free massive gravity},
JHEP \textbf{02} (2012) 126.

\end{thebibliography}

\end{document}
